% Partie : théorie de la démonstration
% Chapitre : calcul des séquents
% Section : exemples de preuves

\documentclass[../../../include/open-logic-section]{subfiles}

\begin{document}

\olfileid[fr]{pt}{seq}{ex}

\olsection{Exemples de preuves}

Pour construire une preuve, il est souvent commode de partir
du séquent final et de remonter vers les axiomes.
On choisit dans le séquent une formule principale,
on écrit au-dessus de lui les prémisses d'une règle
dont il est la conclusion, puis on recommence.
Supposons, par exemple, que l'on veuille prouver
\[(!C \land !D) \lif !E \Sequent !C \lif (!D \lif !E)\]
La formule $!C \lif (!D \lif !E)$ est de la forme
$!A \lif !B$ et figure à droite du séquent.
En comparant le séquent au schéma de la règle
$\RightR{\lif}$ de \Log{G1c},
\[\Axiom$ !A, \Gamma \fCenter \Delta, !B$
\RightLabel{\RightR{\lif}}
\UnaryInf$ \Gamma \fCenter \Delta, !A \lif !B $
\DisplayProof\]
on prend $\Gamma = \{(!C \land !D) \lif !E\}$,
$\Delta = \emptyset$, $!A \ident !C$
et $!B \ident !D \lif !E$.
La dernière inférence d'une preuve pourrait donc être
\[\Axiom$ !C, (!C \land !D) \lif !E \fCenter !D \lif !E$
\RightLabel{\RightR{\lif}}
\UnaryInf$(!C \land !D) \lif !E \fCenter !C \lif (!D \lif !E)$
\DisplayProof\]
En recommençant avec $!D \lif !E$ comme formule principale,
on obtient
\[\Axiom$ !D, !C, (!C \land !D) \lif !E \fCenter !E$
\RightLabel{\RightR{\lif}}
\UnaryInf$!C, (!C \land !D) \lif !E \fCenter !D \lif !E$
\RightLabel{\RightR{\lif}}
\UnaryInf$(!C \land !D) \lif !E \fCenter !C \lif (!D \lif !E)$
\DisplayProof\]
On choisit ensuite $(!C \land !D) \lif !E$ et
on emploie la règle $\LeftR{\lif}$ de \Log{G1c},
\[\Axiom$ \Gamma \fCenter \Delta, !A$
\Axiom$ !B, \Gamma \fCenter \Delta$
\RightLabel{\LeftR{\lif}}
\BinaryInf$ !A \lif !B, \Gamma \fCenter \Delta$
\DisplayProof\]
ce qui donne
\[\Axiom$!D, !C \fCenter !E, !C \land !D$
\Axiom$!D, !C, !E \fCenter !E$
\RightLabel{\LeftR{\lif}}
\BinaryInf$ !D, !C, (!C \land !D) \lif !E \fCenter !E$
\RightLabel{\RightR{\lif}}
\UnaryInf$!C, (!C \land !D) \lif !E \fCenter !D \lif !E$
\RightLabel{\RightR{\lif}}
\UnaryInf$(!C \land !D) \lif !E \fCenter !C \lif (!D \lif !E)$
\DisplayProof\]
En prenant $!C \land !D$ comme formule principale
dans le séquent supérieur gauche et en appliquant
la règle $\RightR{\land}$, on obtient
\[
\Axiom$!D, !C \fCenter !E, !C$
\Axiom$!D, !C \fCenter !E, !D$
\RightLabel{\RightR{\land}}
\BinaryInf$!D, !C \fCenter !E, !C \land !D$
\Axiom$!D, !C, !E \fCenter !E$
\RightLabel{\LeftR{\lif}}
\BinaryInf$ !D, !C, (!C \land !D) \lif !E \fCenter !E$
\RightLabel{\RightR{\lif}}
\UnaryInf$!C, (!C \land !D) \lif !E \fCenter !D \lif !E$
\RightLabel{\RightR{\lif}}
\UnaryInf$(!C \land !D) \lif !E \fCenter !C \lif (!D \lif !E)$
\DisplayProof
\]
Toutes les règles employées jusqu'ici sont communes
à \Log{G1c} et à~\Log{G3c}. Les séquents supérieurs
du fragment ainsi construit ont chacun une formule
commune à leurs deux membres. Ils ne sont toutefois
pas nécessairement des axiomes.

Dans \Log{G1c}, il faut les prouver à partir d'axiomes
de la forme $!A \Sequent !A$. Les règles
d'affaiblissement le permettent. Par exemple,
le séquent supérieur gauche $!D, !C \Sequent !E, !C$
se prouve ainsi :
\[
\Axiom$!C \fCenter !C$
\RightLabel{\LeftR{\Weakening}}
\UnaryInf$!D, !C \fCenter !C$
\RightLabel{\RightR{\Weakening}}
\UnaryInf$!D, !C \fCenter !E, !C$
\DisplayProof
\]
La preuve complète dans \Log{G1c} est donc
\[
\Axiom$!C \fCenter !C$
\RightLabel{\LeftR{\Weakening}}
\UnaryInf$!D, !C \fCenter !C$
\RightLabel{\RightR{\Weakening}}
\UnaryInf$!D, !C \fCenter !E, !C$
\Axiom$!D \fCenter !D$
\RightLabel{\LeftR{\Weakening}}
\UnaryInf$!D, !C \fCenter !D$
\RightLabel{\RightR{\Weakening}}
\UnaryInf$!D, !C \fCenter !E, !D$
\RightLabel{\RightR{\land}}
\BinaryInf$!D, !C \fCenter !E, !C \land !D$
\Axiom$!E \fCenter !E$
\RightLabel{\LeftR{\Weakening}}
\UnaryInf$!C, !E \fCenter !E$
\RightLabel{\LeftR{\Weakening}}
\UnaryInf$!D, !C, !E \fCenter !E$
\RightLabel{\LeftR{\lif}}
\BinaryInf$ !D, !C, (!C \land !D) \lif !E \fCenter !E$
\RightLabel{\RightR{\lif}}
\UnaryInf$!C, (!C \land !D) \lif !E \fCenter !D \lif !E$
\RightLabel{\RightR{\lif}}
\UnaryInf$(!C \land !D) \lif !E \fCenter !C \lif (!D \lif !E)$
\DisplayProof
\]
Sa structure, et en particulier sa hauteur, ne dépend
pas des formules $!C$, $!D$ et~$!E$ choisies.
On peut remplacer ces symboles par des formules
quelconques et obtenir encore une preuve
de~\Log{G1c}. Cette preuve est \emph{schématique}.

Dans \Log{G3c}, les axiomes d'identité ont la forme
$!A, \Gamma \Sequent \Delta, !A$, avec $!A$ atomique.
Le séquent $!D, !C \Sequent !E, !C$ ressemble donc
à un tel axiome, avec $\Gamma = \{!D\}$ et
$\Delta = \{!E\}$, mais il n'en est effectivement
un que si $!C$ est atomique. Le fragment construit
plus haut est une preuve complète de \Log{G3c}
si $!C$, $!D$ et $!E$ sont atomiques.

Pour des formules quelconques, nous n'avons donc pas
encore établi, dans \Log{G3c},
\[\Log{G3c} \Proves (!C \land !D) \lif !E \fCenter !C \lif (!D \lif
!E),\]
En pratique, le fragment suffit néanmoins, à condition
de développer ses feuilles en preuves. En effet,
tout séquent $!A, \Gamma \Sequent \Delta, !A$
se prouve dans \Log{G3c}, même si $!A$ n'est pas
atomique. On le montre par récurrence sur
la profondeur~$\depth{!A}$ de~$!A$.

\begin{defn}\ollabel{defn:depth}
La \emph{profondeur}~$\depth{!A}$ d'une formule
se définit comme suit :
\begin{enumerate}
  \item Si $!A$ est atomique, $\depth{!A} = 0$.
  \item Si $!A \ident \lnot !B$, $!A \ident \lforall[x][!B]$
  ou $!A \ident \lexists[x][!B]$,
  alors $\depth{!A} = \depth{!B} + 1$.
  \item Si $!A \ident (!B \land !C)$,
  $(!B \lor !C)$ ou $(!B \lif !C)$,
  alors $\depth{!A} = \max(\depth{!B},\depth{!C}) + 1$.
\end{enumerate}
\end{defn}

On montre facilement que $\depth{A(x)} = \depth{A(t)}$
pour tout terme~$t$. Dans une règle logique,
la formule principale a une profondeur strictement
supérieure à celle des formules actives issues de
sa décomposition. Une copie de la principale
conservée dans la prémisse garde la même profondeur.
Ce fait ne concerne pas les règles structurelles.
Par exemple :
\begin{align*}
\depth{P(x)} = \depth{P(a)} & = 0,\\
\depth{\lforall[x][P(x)]} & = 1,\\
\depth{(\lforall[x][P(x)] \lif P(a))} & = \max(1,0) + 1 = 2.
\end{align*}
On peut ainsi démontrer des résultats par récurrence
sur~$\depth{!A}$, notamment le suivant.

\begin{prop}\ollabel{prop:G1c-axioms}
Pour toute formule~$!A$, il existe dans $\Log{G1c}$
une preuve de $!A \Sequent !A$ dont tous les axiomes
sont atomiques.
\end{prop}

\begin{proof}
  On raisonne par récurrence sur~$\depth{!A}$.
  Si $\depth{!A} = 0$, alors $!A$ est atomique,
  et $!A \Sequent !A$ constitue déjà une preuve
  de~\Log{G1c} de la forme voulue.
  Si $\depth{!A} > 0$, la formule $!A$ est construite
  à partir de formules de profondeur moindre :
  par exemple $!A \ident \lnot !B$,
  $!A \ident (!B \land !C)$ ou
  $!A \ident \lforall[x][!B(x)]$.
  L'hypothèse de récurrence affirme que, pour toute
  formule $!D$ telle que $\depth{!D} < \depth{!A}$,
  on a $\Log{G1c} \Proves !D \Sequent !D$
  à partir d'axiomes atomiques.

  Si $!A \ident \lnot !B$, alors
  $\depth{!B} < \depth{!A}$. Par hypothèse de
  récurrence, il existe dans \Log{G1c} une preuve~$\pi_1$
  de $!B \Sequent !B$ à partir d'axiomes atomiques.
  On la prolonge en une preuve de
  $\lnot !B \Sequent \lnot !B$ :
  \[
  \AxiomC{}
  \RightLabel{$\pi_1$}
  \Deduce$!B \fCenter !B$
  \RightLabel{\LeftR{\lnot}}
  \UnaryInf$\lnot !B, !B \fCenter $
  \RightLabel{\RightR{\lnot}}
  \UnaryInf$\lnot !B \fCenter \lnot !B$
  \DisplayProof
\]

  Si $!A \ident (!B \land !C)$, alors
  $\depth{!B} < \depth{!A}$ et $\depth{!C} < \depth{!A}$.
  Par hypothèse de récurrence, il existe dans \Log{G1c}
  des preuves~$\pi_1$ et~$\pi_2$ de $!B \Sequent !B$
  et de $!C \Sequent !C$ à partir d'axiomes atomiques.
  On les prolonge en une preuve de
  $!B \land !C \Sequent !B \land !C$ :
  \[
  \AxiomC{}
  \RightLabel{$\pi_1$}
  \Deduce$!B \fCenter !B$
  \RightLabel{\LeftR{\Weakening}}
  \UnaryInf$!B, !C \fCenter !B$
  \RightLabel{\LeftR{\land}}
  \UnaryInf$!B \land !C, !C\fCenter !B$
  \RightLabel{\LeftR{\land}}
  \UnaryInf$!B \land !C, !B \land !C \fCenter !B$
  \AxiomC{}
  \RightLabel{$\pi_2$}
  \Deduce$!C \fCenter !C$
  \RightLabel{\LeftR{\Weakening}}
  \UnaryInf$!B, !C \fCenter !C$
  \RightLabel{\LeftR{\land}}
  \UnaryInf$!B \land !C, !B \fCenter !C$
  \RightLabel{\LeftR{\land}}
  \UnaryInf$!B \land !C, !B \land !C \fCenter !C$
  \RightLabel{\RightR{\land}}
  \BinaryInf$!B \land !C, !B \land !C \fCenter !B \land !C$
  \RightLabel{\LeftR{\Contraction}}
  \UnaryInf$!B \land !C \fCenter !B \land !C$
  \DisplayProof
  \]

  Supposons maintenant $!A \ident \lforall[x][B(x)]$.
  Choisissons une constante~$c$ qui n'apparaît pas
  dans $\lforall[x][B(x)]$. Puisque
  $\depth{B(c)} = \depth{!B(x)} < \depth{!A}$,
  l'hypothèse de récurrence fournit une preuve~$\pi_1$,
  dans \Log{G1c}, de $!B(c) \Sequent !B(c)$
  à partir d'axiomes atomiques. On la prolonge en
  une preuve de
  $\lforall[x][B(x)] \Sequent \lforall[x][B(x)]$ :
  \[
  \AxiomC{}
  \RightLabel{$\pi_1$}
  \Deduce$!B(c) \fCenter !B(c)$
  \RightLabel{\LeftR{\lforall}}
  \UnaryInf$\lforall[x][!B(x)] \fCenter !B(c)$
  \RightLabel{\RightR{\lforall}}
  \UnaryInf$\lforall[x][!B(x)] \fCenter \lforall[x][!B(x)]$
  \DisplayProof
\]

  Les autres cas sont laissés en exercice.
\end{proof}

\begin{prob}
  Compléter la preuve de la \olref{prop:G1c-axioms}
  dans les cas $!A \ident (!B \lor !C)$,
  $!A \ident (!B \lif !C)$ et
  $!A \ident \lexists[x][B(x)]$.
\end{prob}

Dans le cas $!A \ident \lnot !B$, on aurait aussi
pu employer la preuve
\[
\AxiomC{}
\RightLabel{$\pi_1$}
\Deduce$!B \fCenter !B$
\RightLabel{\RightR{\lnot}}
\UnaryInf$\fCenter !B, \lnot !B$
\RightLabel{\LeftR{\lnot}}
\UnaryInf$\lnot !B \fCenter \lnot !B$
\DisplayProof
\]
Le calcul intuitionniste \Log{G1i} autorise au plus
une formule dans le conséquent de chaque séquent.
Dans le premier arbre, \LeftR{\lnot} précède \RightR{\lnot},
et chaque séquent satisfait cette restriction. C'est donc
aussi une preuve dans~\Log{G1i}. Dans le second arbre,
le séquent intermédiaire contient à droite $!B$ et~$\lnot !B$.
%
%
%

Même dans \Log{G1c}, on ne peut pas intervertir
les règles du cas $!A \ident \lforall[x][!B(x)]$.
L'arbre suivant \emph{n'est pas} une preuve :
\[
  \AxiomC{}
  \RightLabel{$\pi_1$}
  \Deduce$!B(c) \fCenter !B(c)$
  \RightLabel{\RightR{\lforall}}
  \UnaryInf$!B(c) \fCenter \lforall[x][!B(x)]$
  \RightLabel{\LeftR{\lforall}}
  \UnaryInf$\lforall[x][!B(x)] \fCenter \lforall[x][!B(x)]$
  \DisplayProof
\]
En effet, la condition sur la constante propre
de~$\RightR{\lforall}$ n'est pas satisfaite :
la constante choisie figure encore dans l'antécédent
de la conclusion de cette inférence.

\begin{prop}\ollabel{prop:G3c-axioms}
Tout séquent de la forme
$!A, \Gamma \Sequent \Delta, !A$ est prouvable
dans~\Log{G3c}, même si $!A$ n'est pas atomique.
\end{prop}

\begin{proof}
La preuve ressemble à celle de la \olref{prop:G1c-axioms},
mais il faut formuler l'hypothèse de récurrence
pour tous les contextes. Le cas
$n = \depth{!A} = 0$ est immédiat :
si $n=0$, $!A$ est atomique, et
$!A, \Gamma \Sequent \Delta, !A$
est un axiome de~\Log{G3c}.

Supposons $n>0$. L'hypothèse de récurrence est la
suivante : pour toute formule $!D$ telle que
$\depth{!D} < n$ et tous contextes $\Pi$ et
$\Lambda$, on a
$\Log{G3c} \Proves !D, \Pi \Sequent \Lambda, !D$.
Considérons le cas $!A \ident (!B \land !C)$.
On applique deux fois l'hypothèse :
avec $!D = !B$, $\Pi = !C, \Gamma$
et $\Lambda = \Delta$, pour obtenir une
preuve~$\pi_1$ de $!B, !C, \Gamma \Sequent \Delta, !B$ ;
puis avec $!D = !C$, $\Pi = !B, \Gamma$
et $\Lambda = \Delta$, pour obtenir une
preuve~$\pi_2$ de $!B, !C, \Gamma \Sequent \Delta, !C$.
Voici la preuve de
$!B \land !C, \Gamma \Sequent \Delta, !B \land !C$ :
\[
\AxiomC{}
\RightLabel{$\pi_1$}
\Deduce$!B, !C, \Gamma \fCenter \Delta, !B$
\RightLabel{\LeftR{\land}}
\UnaryInf$!B \land !C, \Gamma \fCenter \Delta, !B$
\AxiomC{}
\RightLabel{$\pi_2$}
\Deduce$!B, !C, \Gamma \fCenter \Delta, !C$
\RightLabel{\LeftR{\land}}
\UnaryInf$!B \land !C, \Gamma \fCenter \Delta, !C$
\RightLabel{\RightR{\land}}
\BinaryInf$!B \land !C, \Gamma \fCenter \Delta, !B \land !C$
\DisplayProof
\]

Si $!A \ident \lforall[x][!B(x)]$, on choisit
une constante qui n'apparaît ni dans
$\lforall[x][!B(x)]$, ni dans $\Gamma$,
ni dans~$\Delta$. On la note~$c$.
L'hypothèse de récurrence, avec $!D = !B(c)$,
$\Pi = \lforall[x][!B(x)], \Gamma$
et $\Lambda = \Delta$, donne une preuve~$\pi_1$
de $!B(c), \lforall[x][!B(x)], \Gamma \Sequent \Delta, !B(c)$.
On obtient alors
\[
\AxiomC{}
\RightLabel{$\pi_1$}
\Deduce$!B(c), \lforall[x][!B(x)], \Gamma \fCenter \Delta, !B(c)$
\RightLabel{\LeftR{\lforall}}
\UnaryInf$\lforall[x][!B(x)], \Gamma \fCenter \Delta, !B(c)$
\RightLabel{\RightR{\lforall}}
\UnaryInf$\lforall[x][!B(x)], \Gamma \fCenter \Delta, \lforall[x][!B(x)]$
\DisplayProof
\]
La condition sur la constante propre
de $\RightR{\lforall}$ est satisfaite par
le choix de~$c$.
Les cas restants se traitent de même et sont
à compléter dans l'exercice suivant.
\end{proof}

\begin{prob}
  Compléter la preuve de la \olref{prop:G3c-axioms}
  dans les cas $!A \ident (!B \lif !C)$
  et $!A \ident \lexists[x][B(x)]$.
  Traiter également la négation et la disjonction.
\end{prob}

\end{document}
